% page 562 of the facsimile (image p-561 of the scan)
% block 162.6 x 237.7 mm ; left margin 8.0 mm
\pgc{-12.700mm}{0.000mm}{
\l{\cel{3}554}
\l{}
\l{\sou{stimular}. (\sou{trans., ad)}. I. Augmentar la ardoro, la agereso\cel{2}di}
\l{\cel{3}ulu. - II. (\sou{medic.})\cel{2}Augmentar la agado di la funcioni}
\l{\cel{3}organala. - DEFIRS.}
\l{}
\l{\sou{stipo}. (\sou{bot.)}\cel{2}Trunketo di la arbori monokotiledona, sen-rama,}
\l{\cel{3}qua finas per fasko de folii. - L.}
\l{}
\l{\sou{stipendio}. Pensiono gratuita, grantita a lernanto en edukerio}
\l{\cel{3}(liceo, skolo administrata da la guvernistaro, seminario).-}
\l{\cel{3}DEIR.}
\l{}
\l{\sou{stipito}. (\sou{bot.})\cel{2}Parto di la planto, qua departas vertikale}
\l{\cel{3}de la radiko, e portas la folii, la flori e la frukti. -}
\l{\cel{3}EFS.}
\l{}
\l{\sou{stipt}ika. (\sou{medic.}) Di qua la efikiveso e la saporo esas as-}
\l{\cel{3}triktiva. - EFIS.}
\l{}
\l{\sou{stipu}lar. (\sou{trans.}) (\sou{yurocienco}).Enuncar kom kondiciono, en}
\l{\cel{3}kontrato. - DEFIS.}
\l{}
\l{\sou{stivar}. (\sou{trans.})\cel{2}Lokizar la kompozanti di kargajo en la holdo}
\l{\cel{3}di navo. - EIS.}
\l{}
\l{\sou{stofo}. Texuro ek lano, silko, filo, e c., de qua onu facas}
\l{\cel{3}vesti, tapeti, e c.- DEFI.}
\l{}
\l{\sou{stoika}. Qua savas tolerar sua sufro sen plendar - ferme. DEFIRS}
\l{}
\l{\sou{stoko}. (\sou{komerco)}. Quanto de vari qua esas en magazino, sta-}
\l{\cel{3}plo. -\cel{2}(\sou{anke metaf.})\cel{2}- DEF.}
\l{}
\l{\sou{stolo}. (\sou{liturgio katol.}) Larja bendo de stofo kun (kom ornivo)}
\l{\cel{3}tri kruci, quan la sacerdoto pasigas cirkum sua kolo kande}
\l{\cel{3}lu celebras la meso, kande lu donas ad ulu la unciono ex-}
\l{\cel{3}trema. - DEFIS.}
\l{}
\l{\sou{stolono.}\cel{2}(\sou{bot.}) Singla de la radiki qui quaze reptas sur la}
\l{\cel{3}surfaco di sulo, ek la stipiti reptera di ula vejetanti}
\l{\cel{3}(kom ex.: fragiero). -}
\l{}
\l{\sou{stomako}. (\sou{anat}.)Vicero, en la supra regiono di abdomino, parto}
\l{\cel{3}di la tubo nutrivala, posho-forma, en qua la nutrivi trans-}
\l{\cel{3}formesas a chimo. - EFIS.}
\l{}
\l{\sou{stono}. Fragmento de silexo, de quarco, o de roko irgaltra. -DE}
\l{}
\l{\sou{stopar}. (\sou{trans.}) Plenigar (aperturo, truo, lakuno) per enduk-}
\l{\cel{3}tar ulo aden lu.-DE.}
\l{}
\l{\sou{storo}. Peco de mato, o de stofo, movebla adsupre od adinfre, per}
\l{\cel{3}resorto, avan la fenestri di chambro, di veturo, por shir-}
\l{\cel{3}mar kontre la suno-radii. - FS.}
}
